The list-based stress packages (Well, River, General-Head Boundary, and Drain,
among others) share a common structure, so the same guidance applies to all of
them. The per-boundary input values---pumping rates, boundary heads,
conductances, and elevations---are stored in named arrays that are re-read from
the input context at the start of each stress period during Read and Prepare.
These variables are therefore in the \textbf{Period-reset} class: a value set
through the API once before the time loop is overwritten at the next Read and
Prepare and has no effect, whereas a value set every time step (after
\texttt{prepare\_time\_step}) takes effect. The memory address of each variable
is \texttt{MODELNAME/PACKAGENAME/VARNAME}, where \texttt{PACKAGENAME} is the
user-assigned package name (not necessarily the package type).

\subsection{Variables Common to the Stress Packages}

In addition to its value variables, each package's table below includes the
per-boundary variables that are common to all of these packages and are most
often modified through the API: the auxiliary values (\texttt{AUXVAR}), the cell
of each boundary (\texttt{NODELIST}), and the number of active boundaries
(\texttt{NBOUND}). The following variables are also common to these packages but
are not listed in the per-package tables, because they are either fixed or are
recomputed by the program rather than set by the user:

\begin{description}
\item[\texttt{MAXBOUND}] The maximum number of boundaries the package can hold.
This is the allocated size of \texttt{NODELIST}, the value arrays, and
\texttt{AUXVAR}; it is fixed when the package is allocated and is in the
\textbf{Structural / read-only} class. It cannot be increased through the API.
\item[\texttt{HCOF}, \texttt{RHS}] The head-coefficient and right-hand-side
contributions of each boundary, dimensioned by \texttt{MAXBOUND}. These are
recomputed from the boundary input at every formulate step, so they are in the
\textbf{Overwritten} class. They are the mechanism by which a custom API package
injects boundary terms, but they are not the knob for a standard package---set
the boundary input variables instead.
\item[\texttt{SIMVALS}] The simulated flow for each boundary, dimensioned by
\texttt{MAXBOUND}. This is an output, recomputed every time step
(\textbf{Overwritten}).
\end{description}

\subsection{Changing Boundary Locations}

The cells a stress package acts on are not fixed: a boundary can be moved to a
different cell, and boundaries can be activated or deactivated, by editing
\texttt{NODELIST} and \texttt{NBOUND} through the API. Because these variables
are Period-reset, the same timing rule as for the value arrays applies---set them
every time step after \texttt{prepare\_time\_step}, so the change is not
overwritten by Read and Prepare.

\subsubsection{Node numbers in \texttt{NODELIST}}

The entries of \texttt{NODELIST} are one-dimensional node numbers: each is a
single flat index identifying a cell, regardless of the model's discretization
type. The cell identifier supplied in the input---layer-row-column
\texttt{(k, i, j)} for a structured (DIS) grid, layer-cell \texttt{(k, n)} for a
layered unstructured (DISV) grid, or node \texttt{(n)} for a fully unstructured
(DISU) grid---is converted to this single node number, which is what
\texttt{NODELIST} stores; it is not stored in the multi-component form used in the
input.

These node numbers are \emph{reduced} node numbers whenever any cell in the model
has an \texttt{IDOMAIN} value less than one, because cells removed from the active
grid are excluded from the numbering. When no cells are removed, the reduced node
numbers are identical to the user node numbers of the original discretization.

When reduced numbering is in effect, a \texttt{NODELIST} value cannot be
interpreted directly as a position in the original discretization; it must first
be mapped from a reduced node number to a user node number. Use the mapping arrays
\texttt{MODELNAME/DIS/NODEUSER} (reduced node number to user node number) and
\texttt{MODELNAME/DIS/NODEREDUCED} (user node number to reduced node number, zero
for removed cells), substituting \texttt{DISV} or \texttt{DISU} for \texttt{DIS}
according to the grid type. When the model has no removed cells these arrays are
allocated with a length of one, which itself signals that reduced and user node
numbers coincide. The user node number can then be converted to the
\texttt{(k, i, j)}, \texttt{(k, n)}, or \texttt{(n)} cell identifier of the
original grid from the grid dimensions; the same mapping is used in reverse to set
\texttt{NODELIST} from a known cell.

\begin{description}
\item[Moving a boundary] Set \texttt{NODELIST[i]} to the reduced node number of
the target cell. The boundary's value arrays (for example \texttt{COND} and
\texttt{BHEAD}) keep their entry index, so they continue to apply to boundary
\texttt{i} at its new cell.
\item[Activating a boundary] Increase \texttt{NBOUND} (up to \texttt{MAXBOUND})
and fill every array for the new index: \texttt{NODELIST}, the value arrays, and
any \texttt{AUXVAR}. Entries beyond the original \texttt{NBOUND} hold
unspecified memory until they are set, so all of them must be assigned.
\item[Deactivating a boundary] Decrease \texttt{NBOUND}, or move the unwanted
boundary to an inactive cell; boundaries whose cell is inactive contribute
nothing.
\end{description}

\noindent Several constraints and cautions apply. \texttt{MAXBOUND} is the hard
upper limit on the number of boundaries and cannot be grown through the API, so a
model that will gain boundaries at run time must be given a large enough
\texttt{MAXBOUND} in its input (for example, by specifying placeholder
boundaries). The target cell of a moved or activated boundary must be active.
The matrix structure is unaffected, because these packages contribute only to the
diagonal of the boundary cell, which always exists. Finally, anything keyed to a
boundary's position---auxiliary variables, boundary names, observations, and
Mover (MVR) connections---is indexed by the boundary's position in the list and
does not follow the boundary to a new cell, so relocating boundaries can
invalidate those associations.
